\scanpage{195}\mspage{144}
\uline{système inductif} de morphismes dans $\mathrm{Cat}$

\[
f\qquad X_s\xrightarrow{f_s}Y
\]

donne donc un morphisme

\[
\int_{\mathrm{lis}}X_*\xrightarrow{\varphi_{\unsure{Q}}}Y,
\]

où $Q$ désigne le diagramme, de type $S^{\wedge\,\circ}$, formé des $X_s$, de $Y$, des $u^*:X_t\to X_s$ (pour $u:s\to t$ dans $S$) et des $f_s$. De $=$ on aura

\[
\int_{\mathrm{prop}}X_*\xrightarrow{\widetilde\varphi_Q}Y,
\]

et si

\[
\Gamma=(S^\wedge+\text{objet initial}),
\]

on trouve des diagrammes $\widehat Q,\widehat Q^{\sim}$ de type $\Gamma^\circ$ dans $\mathrm{Cat}$. Sauf erreur, les relations entre $\widehat Q,\widehat Q^{\sim},\widehat{Q^\circ},\widehat{Q^\circ}^{\sim}$ sont les $\unsure{{=}^{\mathrm s}|{=}\mathrm s}$ que celles explicitées dans la proposition 14 ci-dessus. Je ne vais pas m'arrêter là-dessus avec \unsure{le soin requis|la sur-rigueur}, sauf en cas de besoin. On dira qu'un diagramme $Q$ définit $Y$ comme $\underset{\unsure{I|1}}{\varinjlim}$ $W$-homotopique, si $\varphi_Q\in W$\deleted{\unsure{,} \illegible{2} \unsure{de} \unsure{$f$}} Il revient au $=$ de dire que $\varphi_{Q^\circ}\in W$, ou que $\widetilde\varphi_Q\in W$, ou que $\widetilde\varphi_{Q^\circ}\in W$\ldots.
% archive sheet number: 195
